\section{Motivating Example and Problem Analysis}

\begin{figure}[htbp]
	\centering
	\includegraphics[width=0.9\textwidth]{figure/motivated_example.png}
	\caption{Motivating Example}
	\label{motivated_example}
\end{figure}

% We use a motivating example to show why rule logic alone is insufficient to determine whether a rule interaction should be retained as a candidate rule conflict. Figure~\ref{motivated_example} shows two zones in a smart home: the bedroom on the left and the greenhouse on the right. The bedroom contains a heater (\circled{1}) and a window (\circled{2}); the greenhouse contains a heater (\circled{3}) and a skylight (\circled{4}). Consider two automation rules configured in the greenhouse: R1 turns on the heater at sunset, and R2 opens the skylight and turns off the heater when the temperature reaches 30\celsius.

% This example reveals three observations. \textbf{User preferences influence whether a rule interaction should be treated as benign or as a candidate rule conflict.} In the greenhouse, the interaction between R1 and R2 can be an expected feedback loop for temperature control. With the same rule logic in the bedroom, however, opening the window at night may create an intrusion risk. In practice, such preferences enter the detection process through the Entity Safety Configuration. \textbf{Physical propagation depends on zone and environmental characteristics.} A heater in one zone does not necessarily affect the temperature in another isolated zone, so a detector that assumes heating always raises temperature may report false positives by inferring non-existent interactions. \textbf{Harmfulness requires explicit safety baselines.} Even when an interaction edge exists, the system still needs preferred safe states to determine whether the result should be treated as a candidate rule conflict. For example, opening the greenhouse skylight may be acceptable, whereas opening the bedroom window may violate the preferred safe states defined in the Entity Safety Configuration.

% These observations motivate our design, which utilizes Physical Channel Configuration to model spatial effects and Entity Safety Configuration to encode user preferences. These inputs enable the system to filter rule interactions into candidate conflicts during static analysis and verify them via dynamic monitoring.

We illustrate the challenges of rule conflict detection through a representative smart home scenario involving two distinct zones. Figure~\ref{motivated_example} depicts a bedroom containing a heater and a window , alongside a greenhouse equipped with a heater and a skylight window. Automation rules are configured to manage the indoor temperature in both areas. For instance, a rule might activate the heater at sunset to maintain warmth. A separate rule triggers the opening of the window or skylight if the temperature exceeds 30\celsius. Although the underlying logic of these rules remains identical across both zones, the resulting physical interactions lead to fundamentally different safety outcomes. Opening the skylight in the greenhouse serves as a benign ventilation mechanism to maintain an optimal growth environment. The identical rule execution in the bedroom creates a security vulnerability by providing a potential entry point for intruders at night.

This scenario highlights several limitations in existing conflict detection approaches. Current tools often assume that any device influencing a common physical channel, such as temperature, constitutes a potential interaction. This assumption ignores the reality that \textbf{physical propagation is highly dependent on spatial boundaries}. A heater operating in the bedroom does not influence the environmental state of the isolated greenhouse. Systems lacking this spatial awareness frequently report non-existent interactions, leading to a high volume of false positives. Furthermore, \textbf{the definition of a rule conflict is related to the specific application scenario and user intent}. As Figure~\ref{motivated_example} demonstrated, the same automation logic can manifest as either a benign environmental adjustment or a severe security vulnerability depending on the context. Purely logic-based detectors fail to distinguish between these cases because they lack the safety baselines required to evaluate the harmfulness of an interaction.

Our approach addresses these issues by decoupling rule interactions from safety violations. We recognize that \textbf{an interaction only qualifies as a conflict if it satisfies both physical reachability and a violation of user-defined safety requirements}. This insight motivates a two-stage analysis detection framework guided by specific environmental and safety configurations. The framework utilizes a physical channel configuration to model the propagation characteristics within each zone. This ensures that only physically possible interactions are mapped into the Rule Interaction Dependency Graph (RIDG). To refine this graph, we incorporate an entity safety configuration that encodes the preferred safe states for each zone. This allows the system to prune benign interactions and focus exclusively on candidates that threaten the security of the environment, which are represented in the Rule Conflict Dependency Graph (RCDG). Finally, the system employs runtime assertion verification to bridge the gap between static modeling and actual execution. By monitoring real-time device states, the framework confirms whether these candidate conflicts are actually triggered, thereby eliminating false alarms stemming from static over-approximation.